% ---
% titre: Problème de proportionnalité - la vente de pâtisserie
% theme: FA
% chapitre: Proportionnalité
% sous_chapitre: Résoudre des problèmes
% niveau: [10e]
% type: EVAL
% tags: [proportionnalite, p-question-TS]
% difficulte: 1
% corrige: true
% ---

\question[2]
Lors de la récréation de mardi, la boulangerie a vendu 23 pains au chocolat. Lundi, elle en avait vendu que 20. Combien de pains au chocolat seront vendus mercredi?

\vspace{10pt}{\footnotesize\raggedleft Espace pour ta démarche et tes calculs\par}
\begin{solutionorgrid}[140pt]
La situation n'est pas proportionnelle.  \textbf{(1.5pt)}
\end{solutionorgrid}

\begin{tcolorbox}[enhanced jigsaw, interior hidden, colframe=box, parbox=false]%
Réponse: 
\sol{\vspace{20pt}}{On ne sait pas, car la situation n'est pas proportionnelle. \textbf{(0.5pt)}}
\end{tcolorbox}



